\subsection{MIRAI}
MIRAI is a fully automatic abstract interpreter for Rust source code, which performs static analysis \cite{facebookexperimental_miraidocumentationoverviewmd_2021,facebookexperimental_facebookexperimentalmirai_2023}. The development has been started in November 2018\footnote{According to the tool's GitHub log \cite{facebookexperimental_first_2018}}.

MIRAI provides four key features:
\begin{itemize}
    \item Locate potential runtime errors without annotations \cite{facebookexperimental_facebookexperimentalmirai_2023}
    \item Verify formal correctness with annotations \cite{facebookexperimental_facebookexperimentalmirai_2023}
    \item Tag analysis with annotations, especially taint analysis \cite{facebookexperimental_miraidocumentationoverviewmd_2021,facebookexperimental_facebookexperimentalmirai_2023}
    \item Generate a call graph \cite{facebookexperimental_miraidocumentationcallgraphmd_2022}
\end{itemize}

The tool's goal is to become widely adopted for Rust static analysis \cite{facebookexperimental_facebookexperimentalmirai_2023}. It promotes itself as a production-ready alternative to research projects like Prusti (see Chapter \ref{cap:prusti}) or Creusot (see Chapter \ref{cap:creusot}), that is efficient and easy to use \cite{facebookexperimental_miraidocumentationoverviewmd_2021}.

It also has a \textit{Request for Proposal} from the Web3 Foundation \cite{web3_foundation_grants-programdocsrfpsopenstatic-analysis-for-runtime-palletsmd_2023}.

\subsubsection{Usability}

MIRAI integrates itself into the Rust compiler, leveraging the generated \acs{mir} for its static analysis \cite{facebookexperimental_miraidocumentationwhypluginmd_2021} and can be called via a cargo wrapper.

The tool can be used without annotating Rust source code for a reachability analysis, to prove that the program doesn't contain a execution path, which causes a runtime error \cite{facebookexperimental_miraidocumentationoverviewmd_2021}. This automatic approach can result in false positives and false negatives, especially when the \acs{mir} is missing for a function, e.g. a function call in another language \cite{facebookexperimental_miraidocumentationoverviewmd_2021}. MIRAI requires an \textit{entry point} for its analysis, such as the main, a test or any public function and can only analyse functions, which are reachable from that point \cite{facebookexperimental_miraidocumentationoverviewmd_2021}. This makes MIRAI an good addition to unit testing, as the coverage of reachable functions should be high.

To use the formal correctness check MIRAI offers, annotations are required in the Rust source code \cite{facebookexperimental_miraidocumentationoverviewmd_2021}. MIRAI supports two crates for importing those: first, the \textit{contracts} crate \cite{karroffel_contracts_2023}, which is more light-weight and accessible to beginners \cite{facebookexperimental_miraidocumentationoverviewmd_2021}, but does not support the more advanced features of MIRAI, like tag analysis \cite{facebookexperimental_miraidocumentationoverviewmd_2021}. These kind of annotations can be compared to Prusti or Creusot annotations, as can bee seen in \autoref{lst:mirai_contracts_annotations_pre_postcondition}. The second crate MIRAI supports is the \textit{mirai-annotations} crate \cite{herman_venter_mirai-annotations_2023}, which provides macros that can be used to annotate the code. This approach doesn't allow constructs like old, as can be seen in \autoref{lst:mirai_mirai_annotations_pre_postcondition}. Is it possible to use both crates simultaneously, so the user can decide which feature from which crate is the most accessible.

\begin{lstlisting}[float, language=Rust, caption=Example Rust code of an increase function with a pre- and postcondition from the contracts crate for MIRAI, label={lst:mirai_contracts_annotations_pre_postcondition}]
#[requires(*val <= u32::MAX - 1)]
#[ensures(*val == old(*val) + 1)]
fn incr(val: &mut u32) {
    *val += 1;
}
\end{lstlisting}

\begin{lstlisting}[float, language=Rust, caption=Example Rust code of an increase function with a pre- and postcondition from the mirai-annotations crate for MIRAI, label={lst:mirai_mirai_annotations_pre_postcondition}]
fn incr(val: &mut u32) {
    precondition!(*val <= u32::MAX - 1);
    let _old = *val;
    *val += 1;
    postcondition!(*val == _old + 1);
}
\end{lstlisting}


\subsubsection{Technology}
%https://github.com/facebookexperimental/MIRAI/blob/main/documentation/Overview.md#abstract-interpretation

%Benutzt Z3 solver under the hood laut quelltext
\subsubsection{Feedback}
\subsubsection{Language Support and Verification Conditions}
%Invariants teilweise ja, quantifiers nein aber kein problem, tag analysis/taint analysis!, diagnostics für library code, call graph generator 
\subsubsection{Development}
MIRAI has been developed by Facebook Research to tackle the lack of a static analysis tool for Rust \cite{facebookexperimental_miraidocumentationoverviewmd_2021} for the development of \textit{Diem} \cite{diem_diemdiem_2023}.

\subsubsection{Activity}
MIRAI is being actively supported with over 1,100 commits in the main branch since 2018 \cite{facebookexperimental_facebookexperimentalmirai_2023}. The lasted release at the time of writing, Version 1.1.8, has been released in May 2023 \cite{facebookexperimental_release_2023} and is considered stable \cite{facebookexperimental_miraidocumentationoverviewmd_2021}. There are currently no plans to further expand the tool, except for bug fixing \cite{facebookexperimental_miraidocumentationoverviewmd_2021}, so it can be considered as feature complete.
% Aktiv, eig feature complete 
\subsubsection{Scientific activity}
% Kein Paper nur GitHub doku
\subsubsection{Documentation}
% Liefert Codebeispiele und gibt doku und dev doku in github (user doku aber eher semi-gut)
